\section{总结与延伸}

\subsection{全课知识图谱}

本课建立了从损失函数到实用训练方法的完整链条：

\begin{center}
\begin{tikzpicture}[node distance=0.8cm, every node/.style={draw, rounded corners, minimum width=3.5cm, minimum height=0.7cm, font=\small}]
\node (loss) {损失函数};
\node (reg) [below=of loss] {正则化};
\node (grad) [below=of reg] {梯度与优化};
\node (opt) [below=of grad] {高级优化器};
\node (lr) [below=of opt] {学习率调度};
\draw[->, thick] (loss) -- (reg);
\draw[->, thick] (reg) -- (grad);
\draw[->, thick] (grad) -- (opt);
\draw[->, thick] (opt) -- (lr);
\node[draw=none, right=1cm of loss, text width=5cm, font=\scriptsize] {Softmax / 交叉熵损失};
\node[draw=none, right=1cm of reg, text width=5cm, font=\scriptsize] {L1, L2, $\lambda$ 超参数};
\node[draw=none, right=1cm of grad, text width=5cm, font=\scriptsize] {数值/解析梯度, 梯度下降, SGD};
\node[draw=none, right=1cm of opt, text width=5cm, font=\scriptsize] {Momentum, RMSProp, Adam, AdamW};
\node[draw=none, right=1cm of lr, text width=5cm, font=\scriptsize] {Step Decay, Cosine, Warmup};
\end{tikzpicture}
\end{center}

\subsection{关键 Takeaways}

\begin{importantbox}{五条核心原则}
\begin{enumerate}
    \item \textbf{正则化是必需的}：仅最小化数据损失必然导致过拟合，正则化通过牺牲训练精度换取泛化能力
    \item \textbf{梯度是优化的基石}：无论使用什么优化器，本质都是在用梯度信息指导参数更新
    \item \textbf{Adam 是默认选择}：除非有特殊原因，否则从 Adam/AdamW 开始
    \item \textbf{学习率需要调度}：固定学习率几乎不是最优选择，预热 + 衰减是标准做法
    \item \textbf{实践胜于理论}：优化器的选择最终是经验性的——尝试几种方案，选效果最好的
\end{enumerate}
\end{importantbox}

\subsection{拓展阅读}

\begin{itemize}
    \item Kingma \& Ba, Adam: A Method for Stochastic Optimization (2015): \url{https://arxiv.org/abs/1412.6980}
    \item Loshchilov \& Hutter, Decoupled Weight Decay Regularization (AdamW, 2019): \url{https://arxiv.org/abs/1711.05101}
    \item Goyal et al., Accurate, Large Minibatch SGD: Training ImageNet in 1 Hour (2017): \url{https://arxiv.org/abs/1706.02677}
    \item CS231N 课程主页: \url{http://cs231n.stanford.edu/}
\end{itemize}

\end{document}
